%-----------------------------------------------------------------------------------------------
\makeatletter
\immediate\write18{datelog > \jobname.info} % site script for $(date '+%Y-%m-%d %Hh%Mm%Ss')
\makeatother
%-----------------------------------------------------------------------------------------------
%-----------------------------------------------------------------------------------------------
\usetheme{Copenhagen}
\usepackage{beamercolorthemeUTF2}
\usefonttheme{serif}
%-----------------------------------------------------------------------------------------------
\usepackage[utf8]{inputenc}
\usepackage[greek,french,english,brazil]{babel} % last becomes the active one
\usepackage{pslatex}
\usepackage{amssymb,amsmath}
\usepackage{soul}
\usepackage[squaren,Gray]{SIunits}
\usepackage[nice]{nicefrac}
\usepackage{tikz}
\usepackage{xspace}
%-----------------------------------------------------------------------------------------------


%-----------------------------------------------------------------------------------------------
%-----------------------------------------------------------------------------------------------
% Mathematical
%-----------------------------------------------------------------------------------------------
\newcommand{\vet}[1]{\underline{{#1}}}
\newcommand{\mat}[1]{\underline{\underline{{#1}}}}
\newcommand{\cub}[1]{\underline{\underline{\underline{{#1}}}}}
\newcommand{\eqdef}{{\ensuremath\stackrel{\text{\tiny def}}{=}}}
%-----------------------------------------------------------------------------------------------
% Presentation
%-----------------------------------------------------------------------------------------------
\newcommand{\BkgImgH}[1]{% Places an image centered on the slide background filling the height
    \usebackgroundtemplate{\parbox{\paperwidth}{%
        \vspace*{1sp}\centering\includegraphics[height=\paperheight]{{#1}}
}}}
\newcommand{\BkgImgW}[1]{% Places an image centered on the slide background filling the width
    \usebackgroundtemplate{\parbox{\paperwidth}{%
        \vspace*{1sp}\centering\includegraphics[width=\paperwidth]{{#1}}
}}}
\newcommand{\ArtEndH}[3]{% Transitions to plain image (last) slide: #1:prefix #2,#3:extensions
    \BkgImgH{root/../art/#1.#2}
    \frame[plain]<handout:0>{%
        \transdissolve\vspace*{80mm}\color{white}\bf\input{root/../art/#1.#3}
}}
\newcommand{\ArtEndW}[3]{% Transitions to plain image (last) slide: #1:prefix #2,#3:extensions
    \BkgImgW{root/../art/#1.#2}
    \frame[plain]<handout:0>{%
        \transdissolve\vspace*{80mm}\color{white}\bf\input{root/../art/#1.#3}
}}
\newcommand{\ImgColW}[3]{% Inserts a full-width image in a column
    \includegraphics[width=\columnwidth]{root/../art/#1.#2}\\[-0.5\baselineskip]
    \parbox{\columnwidth}{\tiny\hfill\scalebox{0.65}{\input{root/../art/#1.#3}}}
}
\newcommand{\txtpic}[1]{%
    \fcolorbox{lightgray}{white!90!black}{{#1}} 
}
%-----------------------------------------------------------------------------------------------


%-----------------------------------------------------------------------------------------------
\title{A.03.02 -- Processos Politrópicos}
\subtitle{(Sistemas Fechados)}
\author{Prof.~C.~Naaktgeboren, PhD}
\date{{\scriptsize\tt%
    \includegraphics[height=6.0mm]{root/00-res/cc/by-nc-nd-88x31.pdf}\\[\smallskipamount]
    https://github.com/CNThermSci/ApplThermSci\\
    Compiled on \input{\jobname.info}
}}
%-----------------------------------------------------------------------------------------------
\begin{document}
%-----------------------------------------------------------------------------------------------
\logo{%
    \parbox{158mm}{% There's a 1mm gap on each side of the 160mm x 90mm slide logo line
        \includegraphics[height=6.5mm]{root/00-res/UTFPR/UTFPR-logo-A.pdf}\hfill%
        \includegraphics[height=6.5mm]{root/00-res/logo/CNThermSci-logo-A.pdf}%
%   \parbox{126mm}{% There's a 1mm gap on each side of the 128mm x 96mm slide logo line
%       \includegraphics[height=6.0mm]{root/00-res/UTFPR/UTFPR-logo-A.pdf}\hfill%
%       \includegraphics[height=6.0mm]{root/00-res/logo/CNThermSci-logo-A.pdf}%
}} % Alpha logos
%-----------------------------------------------------------------------------------------------
\frame{\titlepage}
%-----------------------------------------------------------------------------------------------

%-----------------------------------------------------------------------------------------------
\section{Processos Politrópicos}
%-----------------------------------------------------------------------------------------------

%-----------------------------------------------------------------------------------------------
\subsection{Apresentação}
%-----------------------------------------------------------------------------------------------

    % !j 96 -i8
    %-------------------------------------------------------------------------------------------
    \begin{frame}{Processos Politrópicos -- Definição}\vspace*{-2em}
        \begin{columns}
        \column{0.50\textwidth}
        É todo o processo para o qual:
        $$Pv^n = \mbox{const.}$$
        \uncover<5->{
            A equação é utilizada na forma:
            $$P_1v_1^n = P_2v_2^n$$
        }
        \column{0.50\textwidth}
        \uncover<2->{Onde:}
        \begin{itemize}
            \item<2-> $P$ é a pressão do sistema \\[\bigskipamount]
            \item<3-> $v$ é o volume específico do sistema \\[\bigskipamount]
            \item<4-> $n$ é o \alert{expoente politrópico}
        \end{itemize}
        \end{columns}
    \end{frame}
    %-------------------------------------------------------------------------------------------

    % !j 96 -i8
    %-------------------------------------------------------------------------------------------
    \begin{frame}{Processos Politrópicos -- Apresentação}\vspace*{-2em}
        \begin{itemize}
            \item<1-> Contrário a processos de \alert{propriedade} constante \\[\bigskipamount]
            \item<2-> Um \alert{parâmetro} de processo, \alert{$n$} é mantido
                constante\\[\bigskipamount]
            \item<3-> em processos politrópicos.
        \end{itemize}
    \end{frame}
    %-------------------------------------------------------------------------------------------

    % !j 96 -i8
    %-------------------------------------------------------------------------------------------
    \begin{frame}{Processos Politrópicos -- Apresentação}\vspace*{-2em}
        \begin{align*}
            \uncover<3->{\log\big(}
            Pv^n & \only<1>{= \mbox{const.}}\uncover<2->{= c_1}
            \uncover<3->{\big)\rightharpoondown} \\
            \uncover<4->{
                \log(Pv^n) & = \log(c_1) \equiv c_2 \rightharpoondown \\
            }
            \uncover<5->{
                \log P n\log v & = c_2 \rightharpoondown \\
            }
            \uncover<6->{
                \log P & = c_2 - n\log v \rightharpoondown \\
            }
            \uncover<7->{y & = A + Bx\quad\mbox{for $y\equiv\log{P}$, and $x\equiv\log{v}$...}}
        \end{align*}
    \end{frame}
    %-------------------------------------------------------------------------------------------

%-----------------------------------------------------------------------------------------------
\section{Tópicos de Leitura}
%-----------------------------------------------------------------------------------------------

    %------------------------------------------------------------------------------------------
    \begin{frame}[allowframebreaks]{Tópicos de Leitura}
        \begin{thebibliography}{Çengel, Y.~A., 2013}
            \bibitem[Çengel, Y.~A., 2013]{2013-CengelYA+BolesMA-AMGH}
                Çengel, Y.~A. e Boles, M.~A.
                \newblock{{\em Termodinâmica $7^\mathrm{a}\!$ Edição\/}. \alert{Seção~4-1.}}
                \newblock{\footnotesize AMGH. Porto Alegre. ISBN 978-85-8055-200-3.}
        \end{thebibliography}
    \end{frame}
    %------------------------------------------------------------------------------------------

    % Finishes with stunning image, with credit
    \ArtEndW{amden-875880_1280}{jpg}{txt}

%-----------------------------------------------------------------------------------------------
\end{document}
%-----------------------------------------------------------------------------------------------

